% Компилировать движком XeLaTeX (на Overleaf: Menu -> Compiler -> XeLaTeX).
\documentclass[11pt,a4paper]{article}

\usepackage{fontspec}
\setmainfont{DejaVu Serif}[Scale=0.92]
\usepackage{polyglossia}
\setmainlanguage{russian}
\setotherlanguage{english}
\usepackage[a4paper,margin=2.5cm]{geometry}
\usepackage{amsmath,amssymb}
\usepackage{booktabs}
\usepackage{graphicx}
\usepackage{microtype}
\usepackage[numbers,sort&compress]{natbib}
\usepackage[hidelinks]{hyperref}

\title{\textbf{Интегрированная модель мира: связное описание\\
архитектуры, калибровки и сценариев применения\\[2pt]
\large для задач управления риском}}
\author{Н.~Сотириади}
\date{\today}

\begin{document}
\maketitle

\begin{abstract}
\noindent
В работе описана интегрированная имитационная модель мира, в которой в едином
вычислительном контуре связаны экономика, климат и углеродный цикл, ущерб от
потепления, природные ресурсы, общество и политика, межгосударственные отношения,
культура, а также механизмы кризисов и финансов. Модель построена по принципу
системы взаимодействующих стран (пятьдесят семь страновых агентов и ряд
региональных агрегатов) и продвигается по годам: ежегодно каждый блок обновляет
своё состояние, и результаты передаются в следующий год и в смежные блоки. Цель
работы --- не предъявить точечный прогноз будущего, а представить инструмент,
пригодный для сценарного анализа рисков, ансамблевых расчётов и обнаружения слабых
сигналов приближающихся структурных сдвигов. Изложены проблематика и место модели
среди существующих подходов, её архитектура и лежащая в основе каждого блока
математика и эмпирическая база, процедура калибровки и проверки адекватности с
полученными метриками, явный перечень ограничений, а также типовые сценарии
применения. Экономический и климатический блоки откалиброваны по историческим данным
за 1990--2023~годы и проходят ретроспективную проверку; блоки общества, политики и
геополитики проверяются по более скромному критерию --- превосходству над наивными
эталонными прогнозами.

\medskip
\noindent\textbf{Ключевые слова:} интегрированная модель, управление риском, оценка
неопределённости, сценарный анализ, социальная стоимость углерода, слабые сигналы.
\end{abstract}

\section{Проблематика}

Управление крупными системными рисками --- климатическими, экономическими,
социально-политическими --- требует инструмента, в котором эти области рассматривались
бы совместно, а не порознь. На практике риск-аналитик вынужден опираться на набор
разнородных моделей: климатические модели описывают физику атмосферы, но почти не
содержат финансов и политики; интегрированные модели оценки (например, семейства DICE,
RICE, FUND) соединяют экономику и климат, но, как правило, лишены суверенных финансов,
межгосударственных отношений и социальной динамики \citep{nordhaus2017}; макроэкономические
модели подробно описывают денежно-кредитную сферу, но не учитывают природный капитал.
В результате взаимное влияние сфер --- например, путь от потепления к падению урожаев,
росту цен, социальному напряжению и далее к политической нестабильности и конфликту ---
оказывается за пределами любой отдельной модели.

Между тем именно сопряжение этих каналов порождает наиболее значимые системные риски.
Отсутствие связной модели, охватывающей одновременно климат, экономику, общество и
геополитику, --- это содержательный пробел, ощутимый прежде всего в риск-менеджменте,
где требуется оценивать не отдельный шок, а его распространение по взаимосвязанным
подсистемам. Настоящая работа описывает модель, построенную для заполнения этого пробела.
Подчеркнём с самого начала принципиальную установку: модель не претендует на роль
оракула. Оценки многолетних прогнозов показывают, что даже профессиональные прогнозы
роста на горизонте от двух до пяти лет нередко уступают наивной экстраполяции среднего
\citep{celasun2021}. Поэтому реалистичная цель состоит не в точечном предсказании, а в
\emph{честной откалиброванности относительно собственной неопределённости} и в
способности связно проигрывать сценарии.

\section{Модель и её архитектура}

\subsection{Общее устройство}

Модель представляет мир как совокупность стран, эволюционирующих по годам. Состояние
каждой страны включает экономические, ресурсные, социальные, политические и военные
переменные, а также двусторонние отношения с другими странами; над странами надстроены
глобальные переменные климата и углеродного цикла и мировые цены на ресурсы. Один шаг
модели --- это один год: блоки обновляются в фиксированном порядке, при котором выход
одного блока служит входом следующего, после чего результаты года передаются в
следующий год. Расчётное ядро детерминировано; неопределённость вводится через
ансамблевые прогоны со случайным выбором параметров из обоснованных диапазонов
(раздел~\ref{sec:calib}).

Ниже описаны блоки модели, объединённые в содержательные группы. Для каждого указано,
что он вычисляет, на какой математике и каких данных основан и что он даёт остальной
модели.

\subsection{Экономика}

Выпуск каждой страны определяется тремя факторами: накопленным капиталом, занятым
населением и потребляемой энергией. Производственная функция задана в форме вложенной
функции с постоянной эластичностью замещения (тип \emph{капитал--труд--энергия}),
нормированной так, что в базовой точке калибровки она совпадает с производственной
функцией Кобба--Дугласа \citep{arrow1961}:
\begin{equation}
Y = A \, \Theta \, K^{\alpha} L^{\beta} E^{\gamma},
\qquad \alpha = 0{,}30,\; \beta = 0{,}60,\; \gamma = 0{,}042,
\label{eq:prod}
\end{equation}
где $A$ --- совокупная производительность факторов, $\Theta$ --- технологический
уровень. Вложенная форма с эластичностью замещения капитала и энергии $\sigma_{KE}=0{,}4$
описывает, насколько трудно заместить один фактор другим: при удорожании энергии хозяйство
сдвигается к капиталу и снижает энергопотребление. Именно этот механизм позволяет цене на
углерод реалистично перестраивать экономику. Капитал накапливается по стандартному
правилу
\begin{equation}
K_{t+1} = (1-\delta)\,K_t + s\,Y_t, \qquad \delta = 0{,}05,
\label{eq:capital}
\end{equation}
где норма сбережения $s\approx0{,}24$ модулируется устойчивостью и напряжённостью в
обществе. Производительность $A$ растёт с базовым темпом около одного процента в год,
дополнительно реагируя на расходы на исследования, торговый перелив технологий и
сближение с технологической границей; для прогнозов вперёд её дрейф привязан к
сценарию социально-экономического развития SSP2 \citep{riahi2017}. Рынки энергии,
ресурсов и капитала замыкаются через цены: спрос на энергию и инвестиции реагируют на
цену энергии и стоимость капитала с откалиброванными эластичностями. Отдельно
отслеживаются инфляция (кривая Филлипса с энергетическим переносом издержек),
безработица (закон Оукена), ставка центрального банка (правило Тейлора, \citep{taylor1993}),
бюджет и долг государства, а также торговые потоки. Экономический блок воспроизводит
последнее десятилетие реальной истории и потому является не умозрительным, а
эмпирически закреплённым.

\subsection{Климат и углеродный цикл}

Выпуск сопровождается выбросами углекислого газа; к ним добавляется источник, связанный
с землепользованием. Выбросы накапливаются в атмосфере и распадаются на нескольких
характерных временах, учитываются прочие парниковые газы, а результирующее потепление
рассчитывается стандартной двухрезервуарной моделью теплового баланса поверхности и
океана. Блок откалиброван по мейнстримной научной оценке: он воспроизводит принятую
чувствительность климата к удвоению концентрации углекислого газа (равновесная
чувствительность~$3{,}0$\,$^\circ$C) и отслеживает наблюдаемый ход температуры и
углерода с 1990 по 2023~год \citep{ipcc_ar6_wg1,friedlingstein2023}. Наиболее сильные
самоусиливающиеся обратные связи (таяние вечной мерзлоты, резкий выброс углерода) глубоко
неопределённы и поэтому по умолчанию исключены из основного прогона и исследуются лишь как
хвостовые риски. Это быстрая откалиброванная аппроксимация, уместная в модели экономики и
общества, но не заменяющая полноценную модель земной системы.

\subsection{Ущерб от потепления и стоимость углерода}

Потепление возвращается в экономику как ущерб: повышение температуры снижает выпуск.
Принята квадратичная функция ущерба уровня
\begin{equation}
\Omega(T) = 1 - 0{,}006\,T^{2},
\label{eq:damage}
\end{equation}
что соответствует потерям $2{,}4\,\%$, $5{,}4\,\%$ и $9{,}6\,\%$ валового продукта при
потеплении на $+2$, $+3$ и $+4\,^\circ$C. Эта оценка лежит в середине эмпирического
интервала: примерно в два с половиной раза выше, чем в DICE-2016R2, и ниже центральной
оценки мета-анализа Говарда и Стернера \citep{howard2017,burke2015}. На основе функции
ущерба вычисляется социальная стоимость углерода --- экономический ущерб от выброса
дополнительной тонны углекислого газа, стандартный ориентир климатической политики, ---
методом предельного импульса. Эта величина существенно зависит от схемы дисконтирования
и структуры будущего роста; результаты приведены в разделе~\ref{sec:calib}.

\subsection{Ресурсы}

Для каждой страны отслеживаются запасы, добыча и потребление энергии, продовольствия и
металлов; мировые цены растут, когда спрос превышает предложение, и эта цена передаётся
в инфляцию и далее в экономику. Ресурсный рынок замыкается по цене с откалиброванной
эластичностью спроса; в качестве запасного предусмотрено правило медленной подстройки.
Ресурсы представлены укрупнённо, а не как детальные отраслевые энергосистемы.

\subsection{Общество, политика, геополитика и культура}

Для каждой страны модель отслеживает доверие к власти, социальное напряжение, неравенство
доходов, легитимность институтов и протестное давление. Эти переменные реагируют на
экономику --- рост инфляции и безработицы повышает напряжение и подрывает доверие, а
высокое неравенство усиливает этот эффект --- и, в свою очередь, влияют на устойчивость и
риск кризисов. Каждая пара стран связана отношением: объёмом торговли, наличием союза,
уровнем конфликта или напряжённости, состоянием войны и санкциями. Военный потенциал
каждой страны привязан к наблюдаемым показателям национальной мощи, а не к произвольному
числу \citep{singer1972}. Небольшой набор культурных признаков (прежде всего степень
индивидуализма, избегание неопределённости, долгосрочность ориентации и иерархичность),
а также характер политического режима влияют на экономическое и социальное поведение
\citep{hofstede2010}. В отличие от экономики и климата, для этих связей не существует
единого авторитетного набора данных, поэтому они построены на экспертном суждении и общей
доказательной базе и проверяются по более скромному критерию (раздел~\ref{sec:calib}).
Признаки, которые на деле не влияли на исходы, из модели удалены, чтобы она не
имитировала культурную детальность, которой не использует.

\subsection{Риск, кризисы и финансы}

Поверх плавной годовой динамики модель распознаёт и применяет дискретные кризисы:
долговые, валютные, крах политического режима и климатические экстремальные события, ---
каждый из которых запускается при пересечении соответствующего порога. Тяжесть кризисов
может задаваться распределением с тяжёлым хвостом (большинство событий умеренны, редкие ---
катастрофичны), что отвечает наблюдаемому распределению реальных кризисов и войн
\citep{taleb2007,weitzman2009}; этот режим необязателен и исследуется как
неопределённость. Финансовая сторона ведётся со строгим учётом: всякое изменение долга
уравновешено встречным потоком. Государственные финансы замкнуты, а частный сектор
описан полным банковским балансом, в котором каждый кредит порождает встречный депозит,
так что денежная масса тождественно равна депозитам и кредитам; избыточное заимствование
повышает стоимость кредита и сдерживает экономику --- финансовый акселератор
\citep{bernanke1999,godley2007}.

\subsection{Взаимодействие блоков}

Связность --- определяющее свойство модели. Экономика порождает выбросы; выбросы
определяют потепление; потепление через функцию ущерба~\eqref{eq:damage} снижает выпуск;
цена на углерод удорожает энергию и через замещение в~\eqref{eq:prod} снижает выбросы.
Экономические трудности повышают социальное напряжение и снижают доверие, что меняет
норму сбережения в~\eqref{eq:capital} и повышает риск кризиса; кризис бьёт по выпуску и
капиталу и передаётся партнёрам по торговле. Ресурсный дефицит поднимает цены и инфляцию;
конфликты и санкции нарушают торговлю. Тем самым шок в одной подсистеме распространяется
по остальным --- именно это сопряжение и составляет содержательную ценность модели для
анализа риска.

\section{Калибровка и проверка адекватности}
\label{sec:calib}

\subsection{Источники данных}

Каждый блок привязан к внешним данным (таблица~\ref{tab:data}). Экономический и
климатический блоки откалиброваны и проверены наиболее строго; социально-политические и
геополитические переменные привязаны к признанным индексам, но проверяются по более
мягкому критерию.

\begin{table}[t]
\centering
\caption{Эмпирическая база блоков модели.}
\label{tab:data}
\begin{tabular}{@{}lll@{}}
\toprule
Блок & Источник данных & Назначение \\
\midrule
Экономика & Penn World Table, World Development Indicators & доли факторов, рост, долг \\
Климат и углерод & оценка МГЭИК AR6, Global Carbon Budget & чувствительность, ход $T$ и $\mathrm{CO_2}$ \\
Стоимость углерода & EPA-2023, RFF-SP & ориентир для сверки \\
Прогнозный рост & сценарии SSP2 & дрейф производительности вперёд \\
Конфликты & UCDP Georeferenced Event Dataset & проверка риска конфликта \\
Неравенство & SWIID & неравенство доходов \\
Институты & Worldwide Governance Indicators & легитимность, качество управления \\
Военная мощь & составной индекс национальной мощи & потенциал стран \\
Культура & индексы Хофстеде & культурные признаки \\
\bottomrule
\end{tabular}
\end{table}

\subsection{Принципы калибровки и контроль целостности}

Калибровка следует нескольким правилам. Во-первых, каждый необязательный механизм
выполнен как переключатель, по умолчанию выключенный, если он не закреплён исторической
привязкой; так основное («штатное») поведение модели остаётся привязанным к данным, а
спекулятивные механизмы исследуются отдельно как неопределённость. Во-вторых, для оценки
параметров применяется не точечная подгонка, а \emph{сопоставление с историей}: вместо
поиска единственного «наилучшего» набора параметров отбрасываются области параметрического
пространства, несовместимые с наблюдениями, и остаётся апостериорная область «пока не
отвергнутых» значений --- стандартный подход для дорогих имитаторов \citep{williamson2013}.
В-третьих, климатические коэффициенты, выводимые из исходного состояния, жёстко связаны с
контрольной суммой исходных данных и не могут быть изменены вручную; целостность этой связи
проверяется автоматически. Наконец, действует регрессионный эталон: значимые величины
ретроспективной проверки зафиксированы, и любое изменение ядра должно либо сохранять их с
точностью до допуска, либо сопровождаться явной перекалибровкой.

\subsection{Полученные метрики}

Ретроспективная проверка за 2015--2023~годы (таблица~\ref{tab:backtest}) даёт
среднеквадратические ошибки по валовому продукту, мировым выбросам и температуре. Переход
от прежней производственной функции Кобба--Дугласа к замкнутому ядру с вложенной функцией
замещения и рыночным замыканием улучшил соответствие истории: ошибка по продукту снизилась
с $1{,}03$ до $0{,}59$ триллиона долларов, по выбросам --- с $1{,}61$ до $1{,}15$ гигатонны.

\begin{table}[t]
\centering
\caption{Ретроспективная проверка, 2015--2023.}
\label{tab:backtest}
\begin{tabular}{@{}lc@{}}
\toprule
Показатель & Среднеквадратическая ошибка \\
\midrule
Валовой продукт (20 стран), трлн долл. & $0{,}59$ \\
Мировые выбросы $\mathrm{CO_2}$, Гт & $1{,}15$ \\
Температурная аномалия, $^\circ$C & $0{,}135$ \\
\bottomrule
\end{tabular}
\end{table}

Сопоставление с эталонными прогнозами раскрывает структуру этой точности. Модель
\emph{превосходит} прогноз-постоянство по температуре (оценка мастерства около $+0{,}14$)
и незначительно --- по продукту, но \emph{уступает} наивной линейной экстраполяции по
продукту и выбросам на этом коротком и гладком окне. Это закономерно и согласуется с
оценками профессиональных прогнозов \citep{celasun2021}: структурная модель добавляет
мастерство там, где работает динамика (климат), и проигрывает тривиальной экстраполяции
там, где краткосрочный гладкий агрегат трудно превзойти.

Для блока конфликта применён иной критерий. Показатель конфликтной предрасположенности
стран сопоставлен с историческим реестром вооружённых конфликтов за 1990--2023~годы по
57~странам \citep{sundberg2013}. Модель уверенно ставит реально затронутые конфликтом
страны выше спокойных: площадь под характеристической кривой составляет $0{,}736$, а
оценка мастерства по Брайеру относительно прогноза базовой частоты равна $+0{,}143$
(положительная величина означает наличие мастерства). Это честная планка для подобного
блока: модель пригодна для оценки относительного риска, но не для предсказания конкретных
войн в конкретные даты.

Социальная стоимость углерода чувствительна к схеме дисконтирования. При современной
схеме Рамсея с почти двухпроцентной околонулевой нормой времени она составляет около
$140$ долларов за тонну --- в пределах широкого современного диапазона, чуть ниже
центральных оценок EPA-2023 и RFF-SP (около $190$) \citep{rennert2022,epa2023}; при
дисконтировании по Нордхаусу --- около $43$ долларов. Существенно, что эта величина
\emph{подлинно} чувствительна к структуре будущего роста: по разным вариантам
объективного экономического ядра она колебалась в диапазоне примерно $140$--$380$
долларов. Мы сообщаем это значение вместе с его чувствительностью, а не подгоняем под
целевой ориентир; сама эта чувствительность представляет самостоятельный результат.

\section{Ограничения}

Честный перечень того, чего модель сейчас \emph{не} может, не менее важен, чем перечень
её возможностей.

\textbf{Ожидания обращены в прошлое.} Домохозяйства и фирмы ориентируются на недавнее
прошлое, а не на по-настоящему предвидимое будущее. Поэтому модель не воспроизводит
эффекты упреждающего поведения (например, опережающую реакцию на заранее объявленную
политику).

\textbf{Долгосрочный рост закреплён слабо.} Темп долгосрочного роста производительности
лишь нестрого определён имеющейся короткой историей и для прогнозов вперёд привязан к
внешнему сценарию, а не выведен изнутри модели.

\textbf{Климат --- быстрая аппроксимация.} Климатический блок воспроизводит мейнстримные
показатели, но не является полной моделью земной системы; сильнейшие углеродные обратные
связи глубоко неопределённы и вынесены в хвостовые сценарии.

\textbf{Ущерб широк по всей дисциплине.} Оценки климатического ущерба неопределённы во
всём поле исследований; модель несёт эту неопределённость как интервал, а не как единое
число.

\textbf{Социально-политические связи проверены лишь по знаку.} Для перехода от условий к
социальному недовольству нет единого калибровочного набора данных; эти связи следует
читать как направленно верные и неопределённые, а не как точные прогнозы. Конфликт
по своей природе плохо предсказуем, и блок пригоден лишь для относительного риска.

\textbf{Кризисы основаны на порогах.} Дискретные кризисы задаются пороговыми правилами и
схватывают верное качественное поведение, но не выведены из единого подогнанного набора
данных о кризисах.

\textbf{Передача от денег к ценам --- расширение.} Балансы государства и частного сектора
замкнуты и учётно состоятельны, однако более явная передача от денежной массы к ценовому
давлению остаётся калибруемым расширением.

Общая установка остаётся неизменной: реалистичная цель --- быть откалиброванным
относительно собственной неопределённости, а не служить хрустальным шаром. Честной планкой
точности многолетних прогнозов служит превосходство над наивным эталоном.

\section{Типовые сценарии применения}

Модель предназначена для трёх взаимосвязанных режимов работы, в которых её сопряжённость и
честная оценка неопределённости дают наибольшую отдачу.

\textbf{Сценарный анализ рисков.} Основной режим --- сравнение базовой траектории со
сценарием возмущения: вводится шок или изменение политики (например, цена на углерод,
ресурсный дефицит, обострение отношений между странами), и прослеживается его
распространение по экономике, климату, обществу и геополитике. Содержательно значимы здесь
\emph{относительные} величины --- разности между сценарием и базой, --- а не абсолютные
точечные значения на дальнем горизонте.

\textbf{Ансамблевые расчёты.} Поскольку расчётное ядро детерминировано, неопределённость
вводится прогоном больших ансамблей со случайным выбором параметров из обоснованных
диапазонов. Это даёт не одну траекторию, а распределение исходов, по которому измеряется
чувствительность к каждому параметру и строятся вероятностные оценки. Сценарии
рекомендуется сравнивать именно как ансамбли, а не как одиночные прогоны.

\textbf{Обнаружение слабых сигналов.} Встроенная логика кризисов порогова: кризис
срабатывает при пересечении границы. Обнаружение слабых сигналов --- дополнение к ней:
выявление малозаметного многомерного накопления напряжения, \emph{предшествующего}
структурному сдвигу, до того как сработает любой порог. Для этого применяются три
взаимодополняющих средства. Первое --- расстояние Махаланобиса: для каждого момента
времени оценивается, насколько \emph{совместное} состояние (одновременно продукт,
температура, долг, напряжение и т.\,д.) удалено от его недавнего базового распределения с
учётом корреляций между измерениями; превышение порога $\chi^2$ указывает на переход
системы в необычную область пространства состояний \citep{mahalanobis1936,wilson1931}.
Второе --- обнаружение структурного сдвига: на отдельном ряду ищется смена режима методом
точки разлада со штрафом по информационному критерию, дающая вероятность сдвига и его
наиболее вероятное положение \citep{schwarz1978}. Третье --- индикаторы критического
замедления (рост автокорреляции и дисперсии), нарастающие при приближении к точке
бифуркации \citep{scheffer2009}. Чтобы эти средства выявляли отклонения от \emph{динамики},
а не вторично детектировали плавный тренд, ряды предварительно приводятся к стационарному
приращению. На синтетических данных и на траекториях самой модели средства ведут себя
ожидаемо: внесённый спад продукта на $15\,\%$ распознаётся точно на нужном шаге, тогда как
спокойный прогон остаётся вблизи теоретической частоты ложных срабатываний. Рабочий приём
здесь, как и в сценарном анализе, --- \emph{относительный}: сопоставление профиля слабых
сигналов сценария с базовым (более ранние и более частые отклонения, более высокая
вероятность сдвига) указывает, что сценарий смещает систему к структурному переходу.

\section{Заключение}

Представленная модель связывает в едином контуре экономику, климат, природный капитал,
общество, политику, геополитику и культуру и тем самым заполняет содержательный пробел
между специализированными моделями отдельных областей. Её ценность для управления риском
состоит не в точечном прогнозе, а в способности связно проигрывать сценарии, честно
выражать неопределённость через ансамбли и выявлять слабые сигналы приближающихся
структурных сдвигов. Экономический и климатический блоки эмпирически закреплены и проходят
ретроспективную проверку; социально-политические и геополитические блоки, составляющие
отличительную черту модели, проверены по более скромному, но честному критерию ---
превосходству над наивными эталонами. Дальнейшее развитие связано с введением упреждающих
ожиданий, более глубоким обоснованием долгосрочного роста и явной передачей от денежной
массы к ценам; эти направления оформлены как ограничения и предмет будущей работы, а не как
препятствие к применению модели в очерченных режимах.

\bibliographystyle{plainnat}
\bibliography{references}

\end{document}
